\subsection*{Sprint 4: Creación de la API}
\addcontentsline{toc}{subsection}{Sprint 4: Creación de la API}

El cuarto sprint aborda la construcción de la \textbf{API} de \textbf{Muncher},
el componente que hasta ahora faltaba para que el sistema deje de ser una
interfaz sin sustento y pase a operar sobre información real. La aplicación web
desarrollada en los sprints anteriores resuelve la presentación, pero no dispone
de ningún lugar donde consultar o guardar los datos que muestra; la API será
precisamente ese punto de acceso, el único a través del cual el front-end
obtenga y modifique la información persistente del sistema.

Concentrar el acceso a los datos en un único componente no responde
exclusivamente a una necesidad de organización. Todo aquello que se ejecuta en
el navegador queda expuesto a quien lo utiliza: el código puede inspeccionarse,
las peticiones pueden observarse y modificarse, y ninguna comprobación
efectuada en el cliente puede considerarse fiable, puesto que nada impide
eludirla. Por este motivo, las operaciones cuya ejecución no resulta segura en
el lado del cliente —la validación definitiva de los datos que entran en el
sistema, la aplicación de las reglas de negocio y el manejo de cualquier
credencial de acceso a servicios externos— residirán en la API, donde se
ejecutan en un entorno controlado y no accesible al usuario.

El desarrollo seguirá el mismo criterio aplicado al resto del proyecto:
construir primero en local y desplegar después. La API se incorporará a la
composición de servicios del entorno local, de modo que pueda desarrollarse y
probarse contra el front-end antes de existir en la nube, y solo una vez
comprobado su funcionamiento se llevará al entorno desplegado empleando la
misma definición de infraestructura como código que sostiene los despliegues
actuales.

Finalmente, la API se publicará \textbf{bajo el mismo dominio que el
	front-end}, como una ruta más del punto de entrada ya existente en lugar de
como un servicio con dirección propia. Esta disposición evita las peticiones
entre orígenes distintos y la configuración de CORS que estas exigirían, y
permite que el acceso a la API quede restringido al front-end, que es su único
consumidor previsto: ningún cliente ajeno al sistema deberá poder alcanzarla
directamente.

\subsubsection*{Objetivos iniciales}

\begin{itemize}
	\item Seleccionar el tipo de API: REST, GraphQL u otras alternativas.
	\item Seleccionar la tecnología de la API: lenguaje y framework.
	\item Construir una API con un \textit{endpoint} de prueba.
	\item Seleccionar el mecanismo de comunicación entre el front-end y la API
	      —por ejemplo SWR—, así como la forma de generar los tipos del cliente a
	      partir de la especificación OpenAPI.
	\item Desplegar la API en el entorno local.
	\item Conectar el front-end con la API del entorno local.
	\item Seleccionar la infraestructura sobre la que se desplegará la API.
	\item Desplegar la API en la nube.
\end{itemize}

\subsubsection*{Historias de usuario completadas}

\begin{center}
	\begin{tabularx}{\textwidth}{|l|X|c|}
		\hline
		\textbf{Identificador} & \textbf{Nombre}                              & \textbf{Puntos de historia} \\
		\hline
		\usref{US-CI-06}       & Ejecución local de la API                    & 1                           \\
		\hline
		\usref{US-CI-07}       & Corrección automática del código de back-end & 1                           \\
		\hline
		\usref{US-CI-08}       & Conexión del front-end con el back-end       & 2                           \\
		\hline
		\usref{US-DE-03}       & Despliegue en la nube de la API              & 5                           \\
		\hline
	\end{tabularx}
\end{center}

\subsubsection*{Retrospectiva}

El cuarto sprint concluyó con la API en funcionamiento tanto en el entorno local
como en los dos entornos desplegados, y con el front-end obteniendo de ella los
datos que hasta entonces llevaba escritos en su propio código. El sistema pasó
así a tener la forma que conservará en adelante: una interfaz que no conoce más
origen de información que la API, y una API que es el único componente con
acceso a los datos.

La primera decisión fue el estilo de la API, entre una interfaz REST y una
alternativa basada en GraphQL. Se optó por \textbf{REST} atendiendo a la
naturaleza del consumidor: el front-end es el único cliente previsto y las
operaciones del sistema se corresponden con colecciones de recursos bien
delimitadas —recetas, menús, listas de la compra—, de modo que la capacidad de
GraphQL para que cada cliente componga su propia consulta no resuelve ningún
problema real del proyecto y sí añade una capa de resolución que habría que
construir y mantener. A ello se suma que REST aprovecha directamente la
semántica de HTTP, lo que permite delegar en la infraestructura decisiones como
el tratamiento de la caché en lugar de resolverlas en el código.

Para la implementación se seleccionó \textbf{FastAPI} sobre Python. El factor
determinante no fue el rendimiento del framework, sino que deriva la
especificación OpenAPI de las propias anotaciones de tipos de los
\textit{endpoints}: los modelos que describen las respuestas son a la vez la
validación en ejecución, la documentación publicada y el contrato que consume el
cliente. Esta propiedad se convirtió en el eje de la conexión entre las dos
partes del sistema, que se resolvió generando los tipos de TypeScript a partir
del esquema en lugar de escribirlos a mano. La especificación se trata como un
artefacto intermedio y no se versiona; lo que se guarda en el repositorio es
únicamente el TypeScript resultante, de manera que no puedan existir dos
descripciones divergentes de la misma API. El efecto práctico es que un cambio
en un modelo del back-end que el front-end no haya acompañado deja de compilar,
en vez de manifestarse como un error en ejecución.

El consumo de los datos se apoyó en \textbf{SWR}, que gestiona la petición, su
estado y su revalidación, unido a un cliente tipado construido sobre el esquema
generado. Este cambio obligó a revisar la interfaz de la pantalla de inicio: una
lista de recetas incrustada en el código está siempre disponible, mientras que
una obtenida por red puede tardar, agotar su plazo, no llegar o fallar en el
servidor. Se añadió por ello una clasificación explícita de los modos de fallo
—petición agotada, ausencia de conexión y error del servidor— y estados
diferenciados de carga, lentitud y error, con la posibilidad de reintentar. La
lección es que conectar una interfaz ya maquetada a una fuente remota no es una
sustitución del origen de los datos, sino la incorporación de estados que antes
no existían.

En paralelo se extendieron al back-end los controles de calidad que el front-end
ya tenía. Se incorporaron \textbf{ruff} como formateador y analizador estático y
\textbf{ty} como verificador de tipos, configurados según el mismo criterio
adoptado en el resto del proyecto: se activa el conjunto completo de reglas y
cada exclusión queda escrita junto al motivo que la justifica, de modo que una
regla se desactiva de forma deliberada y no por omisión. Las herramientas se
añadieron a la imagen de contenedor de los controles y se asociaron en
\texttt{pre-commit} a las rutas de la API, aprovechando el encaminamiento por
patrones descrito en el primer sprint, con lo que el \textit{monorepo} admitió
una segunda cadena de herramientas sin que ninguna se aplique a los archivos que
no le corresponden. Durante este sprint se fijó además por escrito la política de
dependencias del proyecto —versión estable más reciente, mantenimiento activo
comprobado y tipos propios—, cuya necesidad se hizo evidente al tener que
resolver en pocos días la elección de varios paquetes en dos ecosistemas
distintos.

El despliegue se abordó buscando el menor coste de operación compatible con los
requerimientos. Se descartó mantener un servidor en ejecución permanente, cuyo
gasto es constante aunque no reciba tráfico, y la API se desplegó como una
\textbf{función Lambda} a la que se accede a través de una URL de función. La
aplicación no se modificó para ello: un punto de entrada específico emplea
\textbf{Mangum} para adaptar los eventos de la función a la interfaz ASGI que el
servidor local invoca, por lo que el código de la API no conoce el entorno en el
que se ejecuta.

La relación entre la plantilla de infraestructura y el código sigue el mismo
patrón ya empleado con el front-end: CloudFormation crea la función con un
paquete de sustitución, porque no puede subir un artefacto que no posee, y el
flujo de despliegue reemplaza después su código, igual que crea el \textit{bucket}
vacío y sube a continuación la aplicación compilada. Esta separación mantiene una
única definición de la infraestructura para ambos entornos y deja el despliegue
del código como un paso independiente y repetible.

La publicación bajo el mismo dominio que el front-end, prevista en los objetivos,
se materializó como un comportamiento adicional de la distribución de CloudFront
que dirige la ruta \texttt{/api} hacia la función. Esta disposición resolvió a la
vez las dos cuestiones que la motivaban: ninguna petición del front-end es de
origen distinto, por lo que no hay configuración de CORS que mantener, y la URL
de la función queda inalcanzable de forma directa, ya que exige autenticación y
solo la distribución del entorno correspondiente está autorizada a firmar sus
peticiones, lo que satisface \reqref{RNF-SE-06}. Como consecuencia añadida, el
muro de autenticación del entorno de desarrollo cubre también la API, de modo que
sus datos no son legibles a través de ella.

Esta misma disposición reveló la dificultad técnica más ilustrativa del sprint.
Servir la API como una ruta del front-end implica que el prefijo \texttt{/api}
forma parte de la dirección pública pero no de las rutas que la aplicación
declara, y cada entorno lo retira en un lugar distinto: en la nube lo absorbe el
adaptador del punto de entrada, mientras que en local es el servidor de
desarrollo el que lo elimina al reenviar la petición al contenedor de la API.

El sprint dejó, en resumen, el sistema completo en su estructura aunque todavía
mínimo en funcionalidad: la colección de recetas que la API devuelve sigue fijada
en el código, sin persistencia detrás. Esa es precisamente la tarea que los
sprints siguientes abordan, y el trabajo realizado aquí la delimita, porque el
contrato entre las dos partes y la vía de despliegue ya están establecidos y
añadir una base de datos no obliga a revisarlos.
